% ===================================================================
%  جدول نمبر 13 — ہوٹل۔ --- کھانے کا ہوٹل۔ قہوہ خانہ۔
%
%  Traduction, faite en 2026, en ourdou d'aujourd'hui, sur l'ido de
%  texto/io. Regles de la colonne : voir 10-jadval-01.tex. Une seule
%  numerotation. Ecrit avec outils/ordo.py sous les yeux. Deux notes
%  d'auteur, toutes deux en gras-italique dans le fac-simile, et le
%  bloc t13-08-1 se poursuit apres la seconde sous la MEME cle.
%
%  LE PERSONNEL DE L'HOTEL EST DEJA NOMME, ET DEPUIS LONGTEMPS. Le
%  tableau met en scene onze metiers de service — femme de chambre,
%  caissiere, portier, chasseur, porteur, blanchisseuse, maitre
%  d'hotel, cuisinier, marmiton, garcon, gerant — et l'ourdou les
%  nomme tous sans emprunter : خادمہ, خزانچی, دربان, چھوٹا, قلی,
%  دھوبن, خانساماں, باورچی, باورچی کا شاگرد, بیرا, منتظم. C'est
%  l'inverse du tableau 10, ou les grades de marine etaient tous
%  anglais. La raison est la meme que pour le chemin de fer : ces
%  metiers-la se sont exerces en ourdou, dans des maisons ou l'on
%  parlait ourdou, bien avant qu'un hotel europeen n'ouvre.
%
%  چھوٹا EST LE PLUS VIF DES ONZE. Le « chasseur » de l'ido est le
%  gamin en livree qu'on envoie faire les courses ; l'ourdou dit
%  simplement چھوٹا, « le petit », et c'est le mot que tout le monde
%  emploie encore dans les hotels et les ateliers du Pendjab. Le mot
%  ne decrit pas la fonction, il decrit l'age — et il porte donc,
%  sans le vouloir, tout ce que le fac-simile de 1926 ne dit pas de
%  ce metier.
%
%  ET نرد, LE TRICTRAC (81), EST PLUS VIEUX QUE LE JEU DE L'IDO. Le
%  mot est persan, il nomme le backgammon depuis les Sassanides, et
%  il est arrive en Europe par l'arabe. La colonne l'ecrit sans
%  hesiter, comme elle ecrit شطرنج pour les echecs (78) et تاش pour
%  les cartes (82). Restent ڈرافٹس et ڈومینو, deux jeux europeens
%  sans nom ourdou : le partage des noms de jeux suit celui des noms
%  de plantes du tableau 8 — ce qui a pousse ici a un nom d'ici.
%
%  DEUX PARADES, ET C'EST LE TIRET QUI TRAVAILLE LE PLUS. Le tableau
%  pose trois fois un genitif a l'envers de l'ourdou : la facture
%  (17) avant le registre (16), les nuages (67) avant l'horizon (66),
%  l'etendard (70) avant l'enseigne (69) et la halle (68). La
%  relative « جو » sert pour les deux premiers cas de tete a
%  complement — le cocher (30) avant son omnibus (29), la malle (31)
%  avant son voyageur (32), le cocher (92) avant sa voiture (88) —,
%  mais un genitif enchaine sur trois termes ne se laisse pas
%  relativiser sans lourdeur : on rejette apres tiret. Le releve des
%  emplois tient maintenant sur trois tableaux et il est net — la
%  relative pour la TETE, le tiret pour la CIRCONSTANCE et pour le
%  genitif a trois etages.
%
%  ET UN TROISIEME MOT DE BROUILLON LAISSE DANS LE TEXTE, DANS LA
%  SEULE COLONNE OURDOUE. « Ordre », seul sur sa ligne, au milieu du
%  bloc t13-02-1 — apres « Wait » au tableau 7 et un autre au tableau
%  9. La colonne tamoule en avait laisse deux. Cinq fois la meme
%  faute, jamais prise par un controle : renvoji.py ne la voit que si
%  le brouillon porte un renvoi, kolonoj.py ne la voyait pas du tout,
%  html.py l'aurait publiee au milieu d'un alinea ourdou.
%
%  kolonoj.py releve desormais toute ligne ENTIEREMENT latine, seule,
%  sans macro ni ponctuation, dans un fichier qui n'est pas latin —
%  et il juge « pas latin » d'apres le fichier lui-meme, pas d'apres
%  une liste de langues : moins d'un dixieme de lettres latines une
%  fois les noms de macro otes. Passe sur les 42 colonnes : ZERO
%  signalement. Le premier essai en rendait zero aussi, mais pour une
%  mauvaise raison — il comptait \VUgras et \textsuperscript parmi
%  les lettres latines, si bien qu'aucun fichier n'etait jamais juge
%  non latin. Un controle qui ne trouve rien doit etre essaye sur une
%  faute fabriquee avant qu'on le croie.
%
%  TRENTE-TROIS PAIRES POUR VINGT-DEUX BLOCS : le rapport le plus
%  eleve du volume. Le tableau ne raconte presque rien, il ATTACHE —
%  le cocher de l'omnibus, la malle du voyageur, la cuisine du
%  sous-sol, l'enseigne de la halle. C'est la confirmation la plus
%  claire de la mesure des tableaux 13, 14 et 15 : ce n'est pas la
%  longueur qui coute, c'est la syntaxe.
% ===================================================================

%%K t13-noto-1 noto
\VUnotes{143.2mm}{%
\textit{\VUgras{(*) کمرے کی تفصیل کے لیے جدول نمبر 6 دیکھیے۔}}%
}

%%K t13-apar-1 apar
\VUsaut{3.0mm}
\VUtitre{15.0pt}{60}{جدول نمبر 13}
\VUsaut{1.6mm}
\VUpk{10.2pt}{40}{ہوٹل۔ --- کھانے کا ہوٹل۔}
\VUsaut{2.0mm}
\VUpk{10.2pt}{40}{قہوہ خانہ۔}
\VUsaut{3.4mm}

%%K t13-01-1 p
1. --- کل شام روآں پہنچ کر میں « \textit{بحرِ اوقیانوس ہوٹل} » میں
ٹھہرا، جس کی سفارش ایک دوست نے کی تھی۔

%%K t13-01-2 p
مجھے ایک خوبصورت \VUgras{کمرہ} \textsuperscript{(2)} دیا گیا —
دوسری \VUgras{منزل} \textsuperscript{(1)} پر —، جہاں میں بہت اچھی طرح
سویا (*)۔

%%K t13-02-1 p
2. --- آج صبح اپنے کمرے سے نکل کر، جہاں \VUgras{خادمہ}
\textsuperscript{(3)} ناشتہ لے آئی تھی، میں \VUgras{تمباکو نوشی کے
کمرے} \textsuperscript{(5)} میں گیا، جو \VUgras{مطالعے کے کمرے}
\textsuperscript{(4)} کے طور پر بھی کام آتا ہے، تاکہ \VUgras{اخبار}
\textsuperscript{(6)} پڑھوں، \VUgras{سگریٹ} \textsuperscript{(8)}
پیتے ہوئے۔ پڑھ چکنے کے بعد میں \VUgras{لفٹ} \textsuperscript{(9)} سے نیچے
اترا اور شہر کی سیر کو نکل گیا۔

%%K t13-03-1 p
3. --- چونکہ میں ہوٹل کے \VUgras{منشی} \textsuperscript{(12)} سے یہ
پوچھنا بھول گیا تھا کہ \VUgras{مشترکہ دستر خوان}
\textsuperscript{(11)} پر کھانا کس وقت لگتا ہے، میں جلد واپس آ گیا
اور اِس کا فائدہ اٹھا کر ہوٹل کے \VUgras{غسل خانے}
\textsuperscript{(10)} میں نہانے چلا گیا۔ پھر میں دوبارہ \VUgras{گلّے}
\textsuperscript{(15)} کی طرف گیا تاکہ اپنے ایک دوست کو ٹیلی فون کروں۔
مگر \VUgras{ٹیلی فون} \textsuperscript{(13)} مصروف تھا ؛ میری پریشانی
دیکھ کر \VUgras{خزانچی خاتون} \textsuperscript{(14)}، جو ایک
\VUgras{بل} \textsuperscript{(17)} درج کر رہی تھیں — \VUgras{مسافروں
کے رجسٹر} \textsuperscript{(16)} میں —، نے \VUgras{دربان}
\textsuperscript{(18)} سے کہا کہ \VUgras{چھوٹے} \textsuperscript{(19)}
کو \VUgras{سائیکل پر} \textsuperscript{(20)} میرا کام کرنے بھیج دے۔

%%K t13-04-1 p
4. --- میں نے اُن کا شکریہ ادا کیا اور اُس \VUgras{قلی}
\textsuperscript{(24)} کو انعام دینے باہر نکلا (وہ محنت کش)، جو
اسٹیشن سے اپنی \VUgras{ہاتھ ریڑھی} \textsuperscript{(25)} پر میرا
\VUgras{ٹوپی کا ڈبّا} \textsuperscript{(26)}، \VUgras{میرا اٹیچی}
\textsuperscript{(27)} اور \VUgras{میرا سفری تھیلا}
\textsuperscript{(28)} لایا تھا۔ \VUgras{ڈاکیا}
\textsuperscript{(21)}، جو ابھی ابھی پہنچا تھا، مجھے ایک \VUgras{خط}
\textsuperscript{(22)} دے گیا۔

%%K t13-05-1 p
5. --- میں کچھ دیر اور دروازے کی دہلیز پر رُکا رہا، \VUgras{خط کے
ڈبّے} \textsuperscript{(23)} کے پاس، تاکہ سڑک کی آمد و رفت دیکھوں۔
\VUgras{کوچوان} \textsuperscript{(30)}، جو \VUgras{اومنی بس}
\textsuperscript{(29)} کا تھا، اپنی گاڑی پر ایک \VUgras{صندوق}
\textsuperscript{(31)} لاد رہا تھا، جو ایک \VUgras{مسافر}
\textsuperscript{(32)} کا تھا ؛ اُس مسافر کو \VUgras{ہوٹل کا مالک}
\textsuperscript{(33)} رخصت کر رہا تھا۔ ایک جوان \VUgras{تار والا}
\textsuperscript{(35)} بغیر جلدی کیے ہوٹل کے ایک گاہک کے لیے
\VUgras{تار} \textsuperscript{(34)} لایا ؛ ایک \VUgras{سیاح}
\textsuperscript{(36)}، کندھے پر اپنا \VUgras{کیمرہ}
\textsuperscript{(38)} لٹکائے، اپنی \VUgras{رہنما کتاب}
\textsuperscript{(37)} دیکھ رہا تھا۔

%%K t13-tit-1 sub
\VUsaut{4.0mm}
\VUpk{10.2pt}{40}{دوپہر کے کھانے کا وقت۔}
\VUsaut{3.0mm}

%%K t13-06-1 p
6. --- جالی کے سامنے ایک بے فکر \VUgras{قاصد} \textsuperscript{(39)}
اپنے \VUgras{بوجھ کے کنڈے} \textsuperscript{(40)} اور اپنے
\VUgras{پالش کے صندوق} \textsuperscript{(41)} کے بیچ اونگھ رہا تھا،
جب کہ ایک جوان \VUgras{دھوبن} \textsuperscript{(42)}، جو کپڑوں کا
\VUgras{ٹوکرا} \textsuperscript{(43)} اٹھائے ہوئے تھی، اُس
\VUgras{خانساماں} \textsuperscript{(44)} کے سنجیدہ چہرے پر ہنس رہی
تھی، جو دروازے کے پاس کھڑا تھا۔

%%K t13-07-1 p
7. --- \VUgras{باورچی} \textsuperscript{(45)} کو دیکھ کر، جو اپنے
\VUgras{باورچی خانے} \textsuperscript{(47)} سے نکلا تھا — وہ باورچی
خانہ \VUgras{تہ خانے} \textsuperscript{(48)} میں ہے — تاکہ
\VUgras{باورچی کے شاگرد} \textsuperscript{(46)} کو کسی کام پر بھیجے،
مجھے یاد آیا کہ کھانے کا وقت قریب ہے، اور میں کھانے کے کمرے میں داخل
ہوا، اُس صحن سے گزر کر جہاں ایک \VUgras{موٹر بان}
\textsuperscript{(50)} اپنی \VUgras{موٹر گاڑی} \textsuperscript{(49)}
کھڑی کر رہا تھا، اور جہاں ایک \VUgras{مستری} \textsuperscript{(52)}
\VUgras{سائیکل سوار} \textsuperscript{(53)} کی ہدایتوں کے مطابق اُس
\VUgras{تیل والی تین پہیوں کی گاڑی} \textsuperscript{(51)} کی مرمت کر
رہا تھا جس کا ابھی حادثہ ہوا تھا۔

%%K t13-08-1 p
8. --- میں نے ایک جوان \VUgras{چھوٹے} \textsuperscript{(54)} سے، جو
\VUgras{بیئر کے گلاس} \textsuperscript{(55)} اٹھائے ہوئے تھا، پوچھا کہ
مشترکہ دستر خوان \VUgras{برآمدے} کے نیچے ہے یا نہیں، اور اُس کے ہاں
کہنے پر میں نے ایک میز پر جگہ لے لی اور اپنے لیے بہت عمدہ کھانا
منگوایا۔

%%K t13-noto-2 noto
\VUnotes{143.2mm}{%
\textit{\VUgras{(*) لغت میں دی ہوئی کھانوں کی فہرست کی مدد سے استاد
نیچے دیا ہوا کھانوں کا نقشہ آسانی سے بدل سکیں گے۔}}%
}

%%K t13-tit-2 sub
\VUsaut{4.0mm}
\VUpk{10.2pt}{40}{عمدہ کھانا۔}
\VUsaut{3.0mm}

%%K t13-09-1 p
9. --- بیرا میرے لیے کھانے کا نقشہ (*) اور شرابوں کی فہرست لے آیا۔
میں نے چن کر پہلے \VUgras{لوازمے} کے طور پر \VUgras{مکھن} کے ساتھ
\VUgras{جھینگے} منگوائے، اور \VUgras{پہلے کھانے} کے طور پر
\VUgras{کان کے انداز کی اوجھڑی۔} میں نے \VUgras{مچھلی} نہیں کھائی،
بلکہ میمنے کی \VUgras{ران} کا ایک حصہ مانگا، اور \VUgras{ترکاری} کے
طور پر بالائی کے ساتھ \VUgras{پالک۔} پھر، چونکہ \VUgras{درمیانی
کھانوں} میں سے کوئی مجھے نہ بھایا، میں نے طرح طرح کی میٹھی چیزیں
منگوائیں — \VUgras{پنیر،} \VUgras{پھل} اور \VUgras{بسکٹ۔} اِن پر میں
نے ایک عمدہ سان-امیلیوں \VUgras{شراب} کا اضافہ کیا، اور اپنے کھانے
سے بہت خوش ہو کر \VUgras{بیرے} \textsuperscript{(57)} کو اچھا
\VUgras{انعام} \textsuperscript{(59)} دے کر میز سے اٹھ گیا۔

%%K t13-10-1 p
10. --- مجھے جانے کو تیار دیکھ کر \VUgras{منتظم} \textsuperscript{(60)}
نے تجویز کیا کہ میں چھجّے پر جا کر \VUgras{بحرِ اوقیانوس}
\textsuperscript{(62)} کا شان دار نظارہ دیکھوں۔ دور
\VUgras{مچھیروں کی چھوٹی کشتیاں} \textsuperscript{(63)}،
\VUgras{بادبانی جہاز} \textsuperscript{(64)} اور ایک بہت بڑا
\VUgras{ڈاک جہاز} \textsuperscript{(65)} دکھائی دینے لگے ؛ اُس جہاز کا
کالا خاکہ سفید \VUgras{بادلوں} \textsuperscript{(67)} پر ابھرتا ہے —
\VUgras{افق} \textsuperscript{(66)} کے بادلوں پر۔ ایک ہلکی ہوا اُس
\VUgras{جھنڈے} \textsuperscript{(70)} کو لہرا رہی تھی جو
\VUgras{نشان} \textsuperscript{(69)} کے اوپر لگا ہے — \VUgras{ہال}
\textsuperscript{(68)} کے نشان کے اوپر۔

%%K t13-tit-3 sub
\VUsaut{3.0mm}
\VUpk{10.2pt}{40}{دل بہلاوے۔}
\VUsaut{3.0mm}

%%K t13-11-1 p
11. --- نظارہ اتنا دل چسپ تھا کہ میں نے طے کیا کہ کل قہوہ خانے کی
چھت پر \VUgras{دوہری دوربین} \textsuperscript{(71)} یا
\VUgras{دوربین} \textsuperscript{(72)} لے کر چڑھوں گا۔

%%K t13-12-1 p
12. --- چھجّا چھوڑ کر میں ہوٹل کے قہوہ خانے کی طرف گیا تاکہ کوئی
\VUgras{مشروب} \textsuperscript{(73)} پیوں اور \VUgras{بلیئرڈ کی بازی}
\textsuperscript{(75)} کھیلوں۔ میں بغیر رُکے اُس قہوہ خانے کے ہال سے
گزرا جہاں بہت سے گاہک \VUgras{شطرنج} \textsuperscript{(78)}،
\VUgras{ڈرافٹس} \textsuperscript{(79)}، \VUgras{ڈومینو}
\textsuperscript{(80)}، \VUgras{نرد} \textsuperscript{(81)} یا
\VUgras{تاش} \textsuperscript{(82)} کھیل رہے تھے، اور سیدھا
\VUgras{بلیئرڈ کے کمرے} \textsuperscript{(74)} میں چڑھ گیا۔

%%K t13-13-1 p
13. --- بازی ختم ہونے پر میں نیچے کی منزل پر اترا اور بیرے سے ایک
\VUgras{لفافہ} \textsuperscript{(84)}، \VUgras{کاغذ}
\textsuperscript{(83)}، \VUgras{ڈاک ٹکٹ} \textsuperscript{(85)}،
\VUgras{قلم} \textsuperscript{(87)} اور \VUgras{سیاہی}
\textsuperscript{(86)} مانگی تاکہ ایک خط لکھ سکوں۔

%%K t13-13-2 p
پھر میں نے ایک \VUgras{کوچوان} \textsuperscript{(92)} کو بلایا، جس کی
\VUgras{بگھی} \textsuperscript{(88)} قہوہ خانے کے سامنے کھڑی تھی۔ میں
نے اُس سے گھنٹے کے حساب سے یا ایک چکر کے کرائے کے بارے میں پوچھا،
اور جب ہماری قیمت طے ہو گئی تو اُس نے میرے لیے \VUgras{گاڑی کا
دروازہ} \textsuperscript{(89)} کھولا اور مجھے بٹھایا۔ وہ پھر اپنی
\VUgras{نشست} \textsuperscript{(91)} پر چڑھا، \VUgras{چابک}
\textsuperscript{(93)} سے گھوڑوں کو تیزی سے چلایا، اور ہم ایک دل کش
سیر پر روانہ ہو گئے۔

\VUsaut{6.0mm}
\VUfilet{22mm}
